% Figure: etch behaviour versus RF bias power in a CCP reactor. The directed
% ion energy rises with applied bias (the self-bias voltage grows), which
% raises the vertical etch rate and the anisotropy but lowers the
% mask-to-film selectivity because the harder ion bombardment also attacks
% the mask. The curves are schematic; the numbers track the CCP case study.
\begin{figure}[htbp]
\centering
\begin{tikzpicture}
\begin{axis}[
  width=0.82\textwidth, height=0.52\textwidth,
  axis y line*=left,
  xlabel={RF bias power $P_{\mathrm{bias}}$ (W)},
  ylabel={etch rate (nm\,min$^{-1}$) and anisotropy ($\times100$)},
  xmin=0, xmax=400, ymin=0, ymax=320,
  grid=both, grid style={enggray!20},
  every axis plot/.append style={thick},
  legend pos=north west, legend cell align=left,
]
  % etch rate: rises and saturates as ion-energy-driven removal dominates
  \addplot[engblue, domain=0:400, samples=80]
    {300*(1-exp(-x/130))};
  \addlegendentry{vertical etch rate}
  % anisotropy (fraction times 100): rises toward 1 with directed ions
  \addplot[engorange, domain=0:400, samples=80]
    {100*(1-exp(-x/70))};
  \addlegendentry{anisotropy $A\times100$}
\end{axis}
\begin{axis}[
  width=0.82\textwidth, height=0.52\textwidth,
  axis y line*=right,
  axis x line=none,
  ylabel={selectivity (film\,:\,mask)},
  xmin=0, xmax=400, ymin=0, ymax=12,
  every axis plot/.append style={thick},
  legend pos=north east, legend cell align=left,
]
  % selectivity: falls as harder ion bombardment erodes the mask
  \addplot[engred, domain=0:400, samples=80, dashed]
    {2+9*exp(-x/120)};
  \addlegendentry{selectivity}
\end{axis}
\end{tikzpicture}
\caption[Etch behaviour versus RF bias power]{Etch behaviour in a
capacitively coupled reactor as the radio-frequency bias power on the wafer
electrode is raised. Higher bias deepens the self-bias voltage and the
sheath, so ions arrive with more directed energy. The vertical etch rate
(left axis, solid blue) climbs and saturates once ion-energy-driven removal
dominates the chemical floor, and the anisotropy (left axis, solid orange,
shown as a fraction times one hundred) approaches one as the directed ion
flux overwhelms the isotropic chemical component. The mask-to-film
selectivity (right axis, dashed red) falls in the same direction, because
the harder bombardment that straightens the side walls also erodes the
mask. The working point of an etch recipe sits where the anisotropy is
adequate and the selectivity has not yet collapsed. The curves are
schematic; the operating numbers are carried in the CCP case study.}
\label{fig:vol04ch13:etch-bias}
\end{figure}
